% =====================================================================
%  FGDH thesis template — minimal starter
%  Forschungsgruppe Digital Health, TU Dresden
%
%  Fill in the places marked « … ». A filling-in guide is in
%  docs/filling-in.md. Compile with pdfLaTeX + BibTeX (latexmk does this
%  for you). Save every file as UTF-8.
%
%  Class options:  hyperref  clickable links/TOC
%                  nat       use the fgdh-thesis-nat (DIN 1505) bibliography style
%                  mp        margin notes
%                  en        switch the template to English (default: German)
%                  xlevel    enable a 4th numbered sectioning level
% =====================================================================
\documentclass[hyperref, nat]{fgdh-thesis}

\begin{document}

% --- Title page -------------------------------------------------------
% Pick the command for your thesis type (all take the same 8 arguments):
%   \bachelortitlepage  \mastertitlepage  \diplomatitlepage  \dissertationtitlepage
%   (\seminartitlepage / \projecttitlepage take 4 — see the documentation)
% Args: {title}{degree}{author}{student-id}{supervisor 1}{supervisor 2}{start date}{end date}
\bachelortitlepage
  {«Titel der Arbeit»}
  {Bachelor of Science}
  {«Vorname Nachname»}
  {«Matrikelnummer»}
  {«Titel und Name der/des Betreuenden»}
  {«Titel und Name der/des Zweitbetreuenden»}
  {«TT.MM.JJJJ»}
  {«TT.MM.JJJJ»}

% --- Front matter -----------------------------------------------------
\begin{preface}
  \abstract
  «Kurze Zusammenfassung der Arbeit (Abstract).»

  \tableofcontents
\end{preface}

% --- Main matter ------------------------------------------------------
\introduction
«Einleitung.»

\section{«Erstes Kapitel»}
\label{sec:erstes-kapitel}
«Text. Verweise mit \ref{...}, Zitate mit \cite{schlüssel}» \cite{beispiel2025}.

% Split large chapters into separate files with \include{dateiname}.

\section{«Zweites Kapitel»}
\label{sec:zweites-kapitel}
«Text.»

% --- Back matter ------------------------------------------------------
\begin{appendix}
  \listoffigures
  \listoftables

  \listofabbreviations
  \abbreviation{z.\,B.}{zum Beispiel}

  % Bibliography: entries live in references.bib; style is fgdh-thesis-nat (DIN 1505).
  \bibliography{references}
  \bibliographystyle{fgdh-thesis-nat}

  \begin{appendices}
    \section{«Anhang»}
    «Ergänzendes Material.»
  \end{appendices}

  % Statutory declaration of authorship.
  \declaration

  % Declaration on the use of AI tools (FGDH). \aideclaration leaves three blank rows
  % to fill in. To disclose tools, pass them in the optional argument — one
  % \aitool{tool}{type and purpose of use}{affected chapter} row each, e.g.:
  %   \aideclaration[%
  %     \aitool{ChatGPT}{Formulierungsvorschläge zur sprachlichen Glättung}{Kapitel~3}%
  %     \aitool{Elicit}{Identifikation relevanter Literatur}{Kapitel~2}%
  %   ]
  \aideclaration
\end{appendix}

\end{document}
